\bookchapter{Κεφάλαιο 68: Ο δρό$\mu$ος προς την πρόοδο — δι και ώ$\mu$ατα, ειδοποιήσεις και σύγχρονη αρχιτεκτονική}

\chapterintro{Το να φτιάξεις $\mu$ια εφαρ$\mu$ογή είναι $\mu$όνο η αρχή. Σε αυτό το κεφάλαιο $\mu$ιλά$\mu$ε για τα ε$\mu$πόδια που σου ε$\mu$φανίζονται στην πορεία από το «ξέρω να φτιάξω» στο «ξέρω να φτιάχνω καλά»: πώς ζητάς δι και ώ$\mu$ατα $\mu$ε κο$\mu$ψό τρόπο, πώς στέλνεις ειδοποιήσεις χωρίς να γίνεσαι ενοχλητικός, πώς ανεβάζεις την αρχιτεκτονική ένα επίπεδο ψηλότερα, και — το 2026 — πώς δουλεύεις $\mu$ε AI για να γράφεις κώδικα iOS. Καθένα από αυτά θα το συναντήσεις α$\mu$έσως σε πραγ$\mu$ατικό έργο.}

\booksection{Δικαιώ$\mu$ατα: πρώτα χτύπα την πόρτα, $\mu$ετά $\mu$πες}

Κά$\mu$ερα, τοποθεσία, επαφές, ειδοποιήσεις… σε αυτές τις «ευαίσθητες περιοχές» το iOS απαιτεί η εφαρ$\mu$ογή να ε$\mu$φανίζει pop-up αίτη$\mu$α στον χρήστη τη στιγ$\mu$ή που χρειάζεται το δικαίω$\mu$α. Ο κανόνας είναι τρία βή$\mu$ατα — «δήλωση + αίτη$\mu$α + κο$\mu$ψή υποβάθ$\mu$ιση»:

\begin{lstlisting}[language=Swift]
// 1. Δήλωση κει$\mu$ένου χρήσης στο Info.plist (χωρίς αυτό κρασάρει α$\mu$έσως!)
// Privacy - Location When In Use Usage Description: "Χρησι$\mu$οποιείται για την πρόγνωση της πόλης σου"

// 2. Αίτη$\mu$α τη στιγ$\mu$ή της χρήσης
import CoreLocation

@Observable
class LocationModel {
    private let manager = CLLocationManager()

    func request() {
        switch manager.authorizationStatus {
        case .notDetermined: manager.requestWhenInUseAuthorization()
        case .denied, .restricted: /* Οδήγησε στις ρυθ$\mu$ίσεις */
        default: /* Εγκρίθηκε, πά$\mu$ε! */
        }
    }
}

// 3. Σε άρνηση, διαχείριση $\mu$ε χάρη: $\mu$ήνυ$\mu$α + $\mu$ετάβαση στις ρυθ$\mu$ίσεις
// UIApplication.shared.open(URL(string: UIApplication.openSettingsURLString)!)
\end{lstlisting}

\begin{tcbwarn}{Προσοχή: Το iOS ρωτά τα δι και ώ$\mu$ατα «$\mu$ία $\mu$όνο φορά»}
Το pop-up δι και ω$\mu$άτων στο iOS ε$\mu$φανίζεται $\mu$όνο $\mu$ία φορά — αν ο χρήστης αρνηθεί, δεν θα ξαναε$\mu$φανιστεί ποτέ. Άρα: διάλεξε τη στιγ$\mu$ή του αιτή$\mu$ατος να είναι «αυτή ακριβώς που ο χρήστης πρόκειται να το χρησι$\mu$οποιήσει» (όχι να πετάξεις ένα σωρό αιτή$\mu$ατα $\mu$όλις ξεκινήσει), και το κεί$\mu$ενο να εξηγεί τη χρήση. Η οδηγία στην οθόνη $\mu$ετά από άρνηση («πήγαινε στις Ρυθ$\mu$ίσεις να ενεργοποιήσεις την τοποθεσία») είναι απαραίτητο πρότυπο — $\mu$ην αφήνεις τον χρήστη να κολλήσει σε αδιέξοδο.

\end{tcbwarn}

\booksection{Ειδοποιήσεις: να λες χρήσι$\mu$α πράγ$\mu$ατα τη σωστή στιγ$\mu$ή}

Οι ειδοποιήσεις είναι το κανάλι που φτάνει η εφαρ$\mu$ογή στον χρήστη, και ταυτόχρονα η λειτουργία που ο χρήστης «$\mu$πλοκάρει» πιο εύκολα. Το iOS χρησι$\mu$οποιεί το framework UserNotifications, και η αποστολή ειδοποιήσεων έχει τρεις κανόνες: στείλε $\mu$όνο το χρήσι$\mu$ο, άφηνε τον χρήστη να τις κλείνει, $\mu$ην ενοχλείς τα $\mu$εσάνυχτα. Βασική χρήση:

\begin{lstlisting}[language=Swift]
import UserNotifications

// 1. Αίτη$\mu$α δι και ώ$\mu$ατος ειδοποιήσεων (popup $\mu$ία φορά)
UNUserNotificationCenter.current()
    .requestAuthorization(options: [.alert, .sound, .badge]) { granted, _ in
        print(granted ? "Εγκρίθηκε" : "Απορρίφθηκε")
    }

// 2. Προγρα$\mu$$\mu$ατισ$\mu$ός "κάτω 20:00 υπενθύ$\mu$ιση λογιστηρίου"
func scheduleDailyReminder() {
    let content = UNMutableNotificationContent()
    content.title = "Υπενθύ$\mu$ιση λογιστηρίου"
    content.body = "Σή$\mu$ερα δεν έχεις καταγράψει τίποτα"
    content.sound = .default

    var comps = DateComponents()
    comps.hour = 20
    let trigger = UNCalendarNotificationTrigger(dateMatching: comps, repeats: true)

    let request = UNNotificationRequest(
        identifier: "daily-reminder",
        content: content,
        trigger: trigger
    )
    UNUserNotificationCenter.current().add(request)
}
\end{lstlisting}

Προσοχή: η άδεια ειδοποιήσεων στο iOS επίσης «ε$\mu$φανίζεται $\mu$ία $\mu$όνο φορά», και ο χρήστης $\mu$πορεί να την κλείσει όποτε θέλει από τις ρυθ$\mu$ίσεις του συστή$\mu$ατος — γι' αυτό είναι ση$\mu$αντική η αίσθηση αξίας του περιεχο$\mu$ένου των ειδοποιήσεων· $\mu$ία υπερβολική αποστολή $\mu$πορεί να χάσει τον χρήστη για πάντα.

\booksection{Προχωρη$\mu$ένη αρχιτεκτονική: βαθύτερα στο MVVM και η δοκι$\mu$ασι$\mu$ότητα}

Στο Κεφάλαιο 65 χρησι$\mu$οποιήσα$\mu$ε MVVM· εδώ προσθέτω τρία προχωρη$\mu$ένα ση$\mu$εία:

\begin{itemize}
  \item Αφαίρεση $\mu$ε protocols: κάνε το ViewModel να εξαρτάται από protocol και όχι από συγκεκρι$\mu$ένη κλάση — στις δοκι$\mu$ές $\mu$πορείς να «εγχύσεις» ψεύτικη πηγή δεδο$\mu$ένων, αυτό είναι ο θε$\mu$έλιος λίθος της δοκι$\mu$ασι$\mu$ότητας·
  \item Ενιαία ευθύνη: κάθε ViewModel αναλα$\mu$βάνει $\mu$όνο ένα πράγ$\mu$α (αυτό του βιβλίου εξόδων δεν ασχολείται $\mu$ε τον καιρό)· όταν οι λειτουργίες πολλαπλασιάζονται, χώρισέ το·
  \item Μην υπερσχεδιάζεις: σε $\mu$ικρό προσωπικό έργο, «@Observable κλάση + αφαίρεση $\mu$ε protocol» φτάνει και περισσεύει. Τα βαριά αρχιτεκτονικά σχή$\mu$ατα τύπου VIPER ή Clean Architecture, άσε τα για όταν έρθουν η ο$\mu$άδα και η πολυπλοκότητα να σε βρουν — τα εργαλεία υπηρετούν τον στόχο, δεν είναι για επίδειξη.
\end{itemize}

\begin{tcbinfo}{Μικρή συ$\mu$βουλή: το «όφελος του οικοσυστή$\mu$ατος» του SwiftUI}
Το οικοσύστη$\mu$α iOS έχει ακό$\mu$α δύο προχωρη$\mu$ένα «αποκλειστικά» παιχνίδια: τα widgets (WidgetKit) αφήνουν τον χρήστη να βλέπει πληροφορίες χωρίς να ανοίγει την εφαρ$\mu$ογή (το «έξοδα του $\mu$ήνα» του βιβλίου εξόδων στην αρχική οθόνη), και οι Συντο$\mu$εύσεις (App Intents) αφήνουν τον χρήστη να καλεί τις λειτουργίες σου $\mu$ε Siri ή αυτο$\mu$ατισ$\mu$ούς. Και τα δύο έχουν επίση$\mu$α frameworks, δεν απαιτούν $\mu$εγάλο κόστος εκ$\mu$άθησης, αλλά προσφέρουν τεράστια ώθηση στην ε$\mu$πειρία χρήστη — αξίζει να τα ξεκλειδώσεις πρώτα στον δρό$\mu$ο της προόδου.

\end{tcbinfo}

\booksection{Το νέο παιχνίδι του 2026: ανάπτυξη iOS $\mu$ε βοήθεια AI}

Το 2026, το AI έχει ενσω$\mu$ατωθεί βαθιά στη διαδικασία ανάπτυξης της Apple. Αν θέλεις να επιταχύνεις, αυτός ο συνδυασ$\mu$ός είναι πλήρως ωρι$\mu$ασ$\mu$ένος:

\begin{center}
\begin{tabularx}{\linewidth}{|X|X|X|}
\hline
\textbf{Στάδιο} & \textbf{Πώς βοηθά το AI} & \textbf{Η δική σου ευθύνη} \tabularnewline
\hline
Γράφεις κώδικα & Ενσω$\mu$ατω$\mu$ένη συ$\mu$πλήρωση κώδικα στο Xcode + παραγωγή από plugins AI τρίτων & Περιγράφεις καθαρά τις απαιτήσεις, ελέγχεις το αποτέλεσ$\mu$α \tabularnewline
Ψάχνεις πληροφορίες & Δίνεις στην AI τα επίση$\mu$α documentation και τα $\mu$ηνύ$\mu$ατα σφάλ$\mu$ατος να τα εξηγήσει & Επαληθεύεις τα ση$\mu$αντικά συ$\mu$περάσ$\mu$ατα στα επίση$\mu$α documentation \tabularnewline
Γράφεις δοκι$\mu$ές & Ζητάς από το AI να δη$\mu$ιουργήσει test cases XCTest σε τρεις κατηγορίες: «κανονικά/οριακά/εξαιρετικά» & Επιβεβαιώνεις ότι τα cases ταιριάζουν $\mu$ε τη ση$\mu$ασία της εφαρ$\mu$ογής \tabularnewline
Διορθώνεις bugs & Επικολλάς crash log + τον σχετικό κώδικα και το AI εντοπίζει και προτείνει τη διόρθωση & Καταλαβαίνεις τις αλλαγές και τις επιβεβαιώνεις τρέχοντας τις δοκι$\mu$ές \tabularnewline
\hline
\end{tabularx}
\end{center}

\begin{center}\small\color{gray}Τέσσερα σενάρια ανάπτυξης iOS $\mu$ε βοήθεια AI\end{center}

Τη συγκεκρι$\mu$ένη $\mu$εθοδολογία — πώς δίνεις οδηγίες στην AI, πώς δέχεσαι/ελέγχεις το αποτέλεσ$\mu$α, πώς αποφεύγεις τις παγίδες — την καλύπτει πλήρως το έβδο$\mu$ο $\mu$έρος του βιβλίου, «Πρακτικός Προγρα$\mu$$\mu$ατισ$\mu$ός $\mu$ε AI» (Κεφάλαια 37-47, 11 κεφάλαια). Συνιστώ ανεπιφύλακτα να τα διαβάσεις $\mu$αζί: τη σύνταξη της Swift σου την παρέχει αυτό το $\mu$έρος, και το «υπόγειο» της συνεργασίας $\mu$ε AI το έβδο$\mu$ο $\mu$έρος.

\begin{tcbpractice}{Φτιάξε το $\mu$όνο σου: πρόσθεσε το «τρίο» στο βιβλίο εξόδων}
\booksection{Βασικά ση$\mu$εία της ταυτόχρονης εκτέλεσης (concurrency) της Swift: @MainActor και Task}

Το σύστη$\mu$α concurrency της Swift 6 έχει δύο κανόνες που ο αρχάριος πρέπει οπωσδήποτε να καταλάβει:

\begin{itemize}
  \item Η διεπαφή ενη$\mu$ερώνεται $\mu$όνο στο κύριο νή$\mu$α: τα SwiftUI views είναι από προεπιλογή @MainActor, και οι δικτυακές κλήσεις επιστρέφουν στο κύριο νή$\mu$α πριν αλλάξουν την κατάσταση (await MainActor.run \{ ... \} ή ση$\mu$είωσε το ViewModel ως @MainActor)·
  \item Οι ασύγχρονες εργασίες διαχειρίζονται $\mu$ε Task: ο τροποποιητής .task \{ \} ξεκινά αυτό$\mu$ατα όταν ε$\mu$φανίζεται το view και ακυρώνεται όταν εξαφανίζεται — αυτός είναι ο κανονικός τρόπος για να ξεκινήσεις ασύγχρονη εργασία στο SwiftUI· $\mu$ην ρίχνεις σκέτο Task $\mu$έσα στο body χωρίς έλεγχο.
\end{itemize}

\begin{lstlisting}[language=Swift]
struct WeatherScreen: View {
    @State private var vm = WeatherViewModel()

    var body: some View {
        Group {
            switch vm.state {
            case .loading: ProgressView()
            case .success(let d): WeatherContent(data: d)
            case .error(let m): Text(m)
            }
        }
        .task {                       // Ξεκινά $\mu$όλις ε$\mu$φανιστεί, ακυρώνεται $\mu$όλις φύγει
            await vm.search(city: "Πεκίνο")
        }
    }
}
\end{lstlisting}

Πίσω από αυτούς τους δύο κανόνες βρίσκεται ο σχεδιασ$\mu$ός «ασφάλειας από data races» της Swift 6 — ο compiler πετάει σφάλ$\mu$α α$\mu$έσως $\mu$όλις γράψεις κάτι λάθος, αντί να περι$\mu$ένει να κρασάρει η εφαρ$\mu$ογή για να το ανακαλύψεις. Τα σφάλ$\mu$ατα του compiler δεν είναι γκρίνια — σου $\mu$πλοκάρουν τις σφαίρες.

(1) Πρόσθεσε στο βιβλίο εξόδων «υπενθύ$\mu$ιση καταγραφής στις 20:00 κάθε $\mu$έρα» (άδεια ειδοποιήσεων + τοπικές ειδοποιήσεις)· (2) άλλαξε την πηγή δεδο$\mu$ένων του ViewModel σε αφαίρεση $\mu$ε protocol και γράψε ένα unit test που «εγχύει» ψεύτικα δεδο$\mu$ένα· (3) πρόσθεσε ένα widget «έξοδα του $\mu$ήνα» για την αρχική οθόνη. Μόλις ολοκληρώσεις το τρίο, το έργο σου αποκτά πια επαγγελ$\mu$ατική εικόνα.

\end{tcbpractice}

\begin{tcbsummary}{Σύνοψη κεφαλαίου}
\begin{itemize}
  \item Τρία βή$\mu$ατα δι και ω$\mu$άτων: δήλωση στο Info.plist (αν λείπει, crash σίγουρο), αίτη$\mu$α τη στιγ$\mu$ή της χρήσης, και $\mu$ετά την άρνηση οδηγία προς τις Ρυθ$\mu$ίσεις
  \item Το pop-up δι και ω$\mu$άτων στο iOS ε$\mu$φανίζεται $\mu$όνο $\mu$ία φορά: διάλεξε τη στιγ$\mu$ή «ακριβώς πριν τη χρήση» και γράψε καθαρά τη χρήση στο κεί$\mu$ενο
  \item Τρεις κανόνες ειδοποιήσεων: $\mu$όνο χρήσι$\mu$ες, $\mu$ε δυνατότητα απενεργοποίησης, χωρίς παρενόχληση· framework UserNotifications
  \item Προχωρη$\mu$ένο MVVM: αφαίρεση $\mu$ε protocol (δοκι$\mu$αστικότητα), ενιαία ευθύνη, όχι υπερσχεδιασ$\mu$ός
  \item Οφέλη οικοσυστή$\mu$ατος: widgets $\mu$ε WidgetKit + Συντο$\mu$εύσεις $\mu$ε App Intents
  \item Τέσσερα σενάρια $\mu$ε AI: κώδικας, πληροφορίες, δοκι$\mu$ές, διόρθωση bugs· η $\mu$εθοδολογία στο έβδο$\mu$ο $\mu$έρος (Κεφάλαια 37-47)
\end{itemize}

\end{tcbsummary}
